% preamble material moved from sample-book.tex for cleanliness and to keep consistent across chapters, which will include this file "remotely"

%% PACKAGES
\usepackage{subfiles}
\usepackage{units}
\usepackage{float}
\usepackage{hyperref}
\usepackage{ulem}
\usepackage{cancel}
\usepackage{tikz}
\usepackage[compat=1.1.0]{tikz-feynman}
\usepackage{feynmf}
\usepackage{imakeidx}
\usepackage{graphicx}
\usepackage{braket}
\usepackage{amsmath}
\usepackage{amsthm}
\usepackage{simplewick}
\usepackage{slashed}
\usepackage{amssymb}
\usepackage[english]{babel}
\usepackage{algorithmic}
\usepackage{floatflt}
\usepackage[ligature, inference]{semantic}
%\usetikzlibrary{external}
\usepackage{stmaryrd}
\usepackage{bm}
%\tikzexternalize[prefix=tikz/]
\usetikzlibrary{quantikz}
\DeclareMathOperator{\Tr}{Tr}
\mathlig{->}{\rightarrow}
\mathlig{|-}{\vdash}
\mathlig{=>}{\Rightarrow}
\mathligson
\usepackage{mathpazo}
\usepackage[mathpazo]{flexisym}
\usepackage{listings}
\lstset{
  basicstyle=\ttfamily,
  mathescape
}
\usepackage{relsize}


\usepackage[framemethod=tikz]{mdframed}
\usepackage[version=4]{mhchem}
\usepackage[thinc]{esdiff} % commands for derivatives
\DeclareMathOperator*{\argmax}{arg\,max}
\DeclareMathOperator*{\argmin}{arg\,min}
\def\inline{\lstinline[basicstyle=\ttfamily, mathescape]}
%\lstMakeShortInline[columns=fixed]|


\theoremstyle{definition}
\newtheorem{example}{Example}[section]
\newtheorem{definition}{Definition}[section]
\newcommand{\ec}{r_1, \ldots r_N}
\newcommand{\nc}{R_1, \ldots R_K}
\newtheorem{conjecture}{Conjecture}[section]

% Make Chapter reference uppercase
\addto\extrasenglish{%
  \renewcommand{\chapterautorefname}{Chapter}%
}
% have chapter numbers...secretly
\usepackage{etoolbox}
\makeatletter
\patchcmd{\Hy@org@chapter}% <cmd>
  {\ifnum}% <search>
  {\refstepcounter{chapter}\ifnum}% <replace>
  {}{}% <success><failure>
\makeatother

% equation/figure numbering
\usepackage{chngcntr}
\counterwithin{figure}{chapter}
\counterwithin{equation}{chapter}

% exclude ToC, LoF, LoT from ToC but include LoS, LoA
\usepackage[nottoc]{tocbibind} 

% to make lists of symbols/acronyms
% remove intoc option if not wanted in TOC
\usepackage[intoc]{nomencl}
\renewcommand{\nomname}{List of Symbols}
% nonumberlist = don't print page number of definition of term in acronym list
\usepackage[acronym, nonumberlist]{glossaries}
\glstoctrue % comment out if not wanted in TOC

\usepackage{booktabs}

% The fancyvrb package lets us customize the formatting of verbatim
% environments.  We use a slightly smaller font.
\usepackage{fancyvrb}
\usepackage{xspace}
\usepackage{cleveref} % have to load this last

% If they're installed, use Bergamo and Chantilly from www.fontsite.com.
% They're clones of Bembo and Gill Sans, respectively.
%\IfFileExists{bergamo.sty}{\usepackage[osf]{bergamo}}{}% Bembo
%\IfFileExists{chantill.sty}{\usepackage{chantill}}{}% Gill Sans

%% SOME GENERAL CONFIGURATION STUFF

\setkeys{Gin}{width=\linewidth,totalheight=\textheight,keepaspectratio}
\graphicspath{{figures/}}
\fvset{fontsize=\normalsize}
\hypersetup{colorlinks}% uncomment this line if you prefer colored hyperlinks (e.g., for onscreen viewing)

%% CUSTOM COMMANDS

% Prints argument within hanging parentheses (i.e., parentheses that take
% up no horizontal space).  Useful in tabular environments.
\newcommand{\hangp}[1]{\makebox[0pt][r]{(}#1\makebox[0pt][l]{)}}

% Prints an asterisk that takes up no horizontal space.
% Useful in tabular environments.
\newcommand{\hangstar}{\makebox[0pt][l]{*}}

% Prints the month name (e.g., January) and the year (e.g., 2008)
\newcommand{\monthyear}{%
  \ifcase\month\or January\or February\or March\or April\or May\or June\or
  July\or August\or September\or October\or November\or
  December\fi\space\number\year
}

% Prints an epigraph and speaker in sans serif, all-caps type.
\newcommand{\openepigraph}[2]{%
  %\sffamily\fontsize{14}{16}\selectfont
  \begin{fullwidth}
  \sffamily\large
  \begin{doublespace}
  \noindent\allcaps{#1}\\% epigraph
  \noindent\allcaps{#2}% author
  \end{doublespace}
  \end{fullwidth}
}

% Inserts a blank page
\newcommand{\blankpage}{\newpage\hbox{}\thispagestyle{empty}\newpage}

% Typesets the font size, leading, and measure in the form of 10/12x26 pc.
%\newcommand{\measure}[3]{#1/#2$\times$\unit[#3]{pc}}

\newcommand{\hlred}[1]{\textcolor{Maroon}{#1}}% prints in red
\newcommand{\hangleft}[1]{\makebox[0pt][r]{#1}}
\newcommand{\hairsp}{\hspace{1pt}}% hair space
\newcommand{\hquad}{\hskip0.5em\relax}% half quad space
\newcommand{\TODO}{\textcolor{red}{\bf TODO!}\xspace}
\newcommand{\na}{\quad--}% used in tables for N/A cells


% sub- and supertext: if sub-/superscripts aren't variables, they shouldn't be italicized
\newcommand{\sub}[1]{_{\text{#1}}}
\renewcommand{\sup}[1]{^{\text{#1}}} %bad practice but I don't think there will be any supremum-taking in the book...

% some particular cases that come up a lot...
\newcommand{\kB}{k_{\text{B}}}
\newcommand{\subred}{_{\text{red}}}
\newcommand{\subox}{_{\text{ox}}}
\newcommand{\supac}{^{\text{ac}}}
\newcommand{\subeq}{_{\text{eq}}}
\newcommand{\substd}{_{\text{std}}}
\newcommand{\subA}{_{\text{A}}}
\newcommand{\subB}{_{\text{B}}}
\newcommand{\subAB}{_{\text{AB}}}
\newcommand{\submix}{_{\text{mix}}}